\documentclass[
  reprint,
  superscriptaddress,
  amsmath,
  amssymb,
  aps,
  pre,
  longbibliography
]{revtex4-2}

% Keep \bm lighter by reserving fewer bold math alphabets.
\newcommand\hmmax{0}
\newcommand\bmmax{0}

% Text and math fonts.
\usepackage[utf8]{inputenc}
\usepackage[lining,semibold]{libertine}
\usepackage[libertine,cmintegrals,bigdelims,vvarbb]{newtxmath}

% Core math support.
\usepackage{amsmath}
\usepackage{amsfonts}
\usepackage{bm}

% Figures and tables.
\usepackage{graphicx}
% REVTeX's table patches require the legacy array implementation on TeX Live 2025.
% Load it before dcolumn, which otherwise loads the incompatible current array.
\usepackage{array}[=2016-10-06]
\usepackage{dcolumn}
\usepackage{booktabs}
\usepackage{siunitx}

% Links and references. Load hyperref before cleveref.
\usepackage{xcolor}
\definecolor{webgreen}{rgb}{0,.5,0}
\definecolor{webbrown}{rgb}{.6,0,0}
\definecolor{RoyalBlue}{rgb}{0.0,0.14,0.4}

\usepackage{hyperref}
\hypersetup{
  colorlinks=true,
  linktocpage=true,
  breaklinks=true,
  pdfstartpage=1,
  pdfstartview=Fit,
  pdfpagemode=UseOutlines,
  plainpages=false,
  bookmarksnumbered=true,
  bookmarksopen=true,
  bookmarksopenlevel=1,
  hypertexnames=true,
  pdfhighlight=/O,
  urlcolor=webbrown,
  linkcolor=RoyalBlue,
  filecolor=RoyalBlue,
  citecolor=webgreen,
  pdftitle={Mutual linearity from classical Markov response theory},
  pdfauthor={Timur Aslyamov and Massimiliano Esposito},
  pdfcreator={pdfLaTeX},
  pdfproducer={LaTeX REVTeX}
}

\usepackage[capitalize]{cleveref}
\displaywidowpenalty=10000
\AtBeginEnvironment{thebibliography}{\interlinepenalty=10000}
\newcommand{\creflastconjunction}{, and\nobreakspace}
\newcommand{\crefrangeconjunction}{--}

% Prevent math macros from leaking into PDF bookmarks.
\pdfstringdefDisableCommands{%
  \def\theta{}%
  \def\varphi{}%
  \def\mathcal{}%
  \def\mathbb{}%
}

% Normal-size equation numbers.
\makeatletter
\def\maketag@@@#1{\hbox{\m@th\normalfont\normalsize#1}}
\makeatother

% Small reusable macros.
\newcommand{\me}{\mathrm{e}}
\newcommand{\tr}{\operatorname{tr}}
\newcommand{\mb}[1]{\mathbf{#1}}
\newcommand{\E}{\mathbb{E}}
\newcommand{\Prob}{\mathbb{P}}
\newcommand{\dd}{\mathrm{d}}
\newcommand{\src}{\operatorname{s}}
\newcommand{\tgt}{\operatorname{t}}


\begin{document}
%%%%%%%%%%%%%%%%%%%%%%%%%%%%%%%%%%%%% frontmatter %%%%%%%%%%%%%%%%%%%%%%%%%%%%%%%%%%%%%
\title{Mutual linearity from classical Markov response theory}

\author{Timur Aslyamov}
\email{timur.aslyamov@uni.lu}
\affiliation{Complex Systems and Statistical Mechanics, Department of Physics and Materials Science, University of Luxembourg, 30 Avenue des Hauts-Fourneaux, L-4362 Esch-sur-Alzette, Luxembourg}

\author{Massimiliano Esposito}
\email{massimiliano.esposito@uni.lu}
\affiliation{Complex Systems and Statistical Mechanics, Department of Physics and Materials Science, University of Luxembourg, 30 Avenue des Hauts-Fourneaux, L-4362 Esch-sur-Alzette, Luxembourg}

\date{\today}


\begin{abstract}
Mutual linearity of stationary probabilities and currents in Markov jump
processes follows directly from classical Markov response theory.
Using the exact finite-response identity expressed through the Drazin
inverse, we show that changing the rates of one edge changes all
stationary probabilities by the same scalar factor multiplied by fixed
coefficients. Eliminating this factor gives the known mutual-linearity
relations, valid arbitrarily far from equilibrium.
\end{abstract}
\maketitle

\section{Introduction}
\label{sec:introduction}

Stationary probabilities and currents generally depend nonlinearly on
transition rates. Nevertheless, varying the rates of one edge produces
affine relations between these quantities, known as mutual
linearity~\cite{Harunari2024,Khodabandehlou2025,BebonSpeck2026,ZhengLu2026,VoitsSchwarz2026}.
Our goal is to give a short common derivation from the classical
finite-response identity of Markov-chain theory~\cite{Schweitzer1968,Meyer1980}.
Written with the Drazin inverse, this identity shows that all probability
changes share one scalar factor. Eliminating it yields the known
relations and makes their extension beyond Markov processes transparent.

\section{Exact stationary response}
\label{sec:response}

The finite-response identity below dates back to the
1960s~\cite{Schweitzer1968}. We simply recast it here for Markov jump
processes.

We consider a finite irreducible Markov jump process,
$\partial_t\bm p=W\bm p$, where $W_{ij}$ is the transition rate from
$j$ to $i$, and $W_{jj}=-\sum_{i\ne j}W_{ij}$.
The stationary distribution satisfies $W\bm\pi=0$ and
$\bm1^\intercal\bm\pi=1$.
For a second irreducible generator $W'=W+\Delta W$, the classical
fundamental-matrix identity takes the form
\begin{align}
    \Delta\bm\pi\equiv\bm\pi'-\bm\pi=-G\Delta W\bm\pi'\,,
    \label{eq:finite}
\end{align}
where $G=W^D$ is the Drazin inverse~\cite{Meyer1975,crook2018drazin,aslyamov2025excess}.
It satisfies $GW=WG=I-\bm\pi\bm1^\intercal$,
$G\bm\pi=0$, and $\bm1^\intercal G=0$.
Indeed, subtracting the two stationary equations gives
$W\Delta\bm\pi=-\Delta W\bm\pi'$.
Multiplication by $G$ and normalization yield \cref{eq:finite},
without assuming that $\Delta W$ is small.

\section{Mutual linearity}
\label{sec:mutual-linearity}

\subsection{Probabilities and observables}

We change $W_{ij}$ and $W_{ji}$,
keeping all other off-diagonal rates fixed and adjusting the diagonal
entries to preserve probability conservation.
Only components $i$ and $j$ of $\Delta W\bm\pi'$ are nonzero,
and they have opposite signs. Equation~\eqref{eq:finite} therefore gives
\begin{align}
    \Delta\pi_n
    &=(G_{nj}-G_{ni})
      (\Delta W_{ij}\pi_j'-\Delta W_{ji}\pi_i')
    \nonumber\\
    &\equiv h_n\alpha.
    \label{eq:edge-response}
\end{align}
Here $h_n\equiv G_{nj}-G_{ni}$ is calculated from the rates before the change,
including $W_{ij}$ and $W_{ji}$.
We keep this initial $W$ fixed and vary $W_{ij}'$ and $W_{ji}'$;
all dependence on these new rates enters
$\alpha \equiv \Delta W_{ij}\pi_j'-\Delta W_{ji}\pi_i'$.
Repeating \cref{eq:edge-response} for state $m$ and eliminating $\alpha$ gives
\begin{align}
    \Delta\pi_n=\frac{h_n}{h_m}\Delta\pi_m\,,
    \qquad h_m\ne0\,.
    \label{eq:probabilities}
\end{align}
Thus $(\pi_m',\pi_n')$ lies on one straight line for every choice of
the new edge rates~\cite{BebonSpeck2026}.
Choosing another starting point changes the $h_n$ but leaves their ratios,
and hence this line, unchanged.

For state observables $O=\sum_n O_n\pi_n$ and
$Q=\sum_n Q_n\pi_n$ with fixed weights, the same argument gives
\begin{align}
    \Delta O=\frac{\sum_n O_n h_n}{\sum_n Q_n h_n}\Delta Q\,,
    \label{eq:observables}
\end{align}
provided the denominator is nonzero.
This also covers mean jump-counting rates on unchanged edges, since
they are fixed linear combinations of the stationary probabilities.
Taking infinitesimal changes in \cref{eq:edge-response} further yields
\begin{align}
    \partial_{W_{ij}}\pi_n=h_n\pi_j\,,\qquad
    \partial_{W_{ji}}\pi_n=-h_n\pi_i\,,
    \label{eq:derivatives}
\end{align}
so that $-\partial_{W_{ij}}\pi_n/\partial_{W_{ji}}\pi_n=\pi_j/\pi_i$
whenever $h_n\ne0$, recovering the response ratio of
Ref.~\cite{BebonSpeck2026}.
Our dynamical fluctuation-response theory gives the same ratio for
responses to impulsive rate perturbations at time $t$, with
$\pi_j/\pi_i$ replaced by $p_j(t)/p_i(t)$~\cite{AslyamovEsposito2026}.
This relation does not require mutual linearity under finite rate changes.

\subsection{Currents}

The stationary current from $l$ to $k$ is
$J_{kl}=W_{kl}\pi_l-W_{lk}\pi_k$.
For an unchanged edge, \cref{eq:edge-response} gives
\begin{align}
    \Delta J_{kl}=(W_{kl}h_l-W_{lk}h_k)\alpha\,.
    \label{eq:output-current}
\end{align}
On the perturbed edge, the rates themselves also change:
\begin{align}
    \Delta J_{ij}
    &=W_{ij}\Delta\pi_j-W_{ji}\Delta\pi_i
      +\Delta W_{ij}\pi_j'-\Delta W_{ji}\pi_i'
    \nonumber\\
    &=(1+W_{ij}h_j-W_{ji}h_i)\alpha\,.
    \label{eq:input-current}
\end{align}
Eliminating $\alpha$ between the currents on $(i,j)$ and another edge
$(k,l)$ gives
\begin{align}
    \Delta J_{kl}
    =\frac{W_{kl}h_l-W_{lk}h_k}
          {1+W_{ij}h_j-W_{ji}h_i}\Delta J_{ij}\,,
    \label{eq:currents}
\end{align}
which is the mutual linearity of stationary
currents~\cite{Harunari2024,Khodabandehlou2025}.
Both current changes contain the same $\alpha$, so their ratio is
independent of the new edge rates, provided the denominator is nonzero
and the perturbed process remains irreducible.

The connection with the first-passage formulation of
Ref.~\cite{Khodabandehlou2025} follows from
$h_n=\pi_n(\tau_{j\to n}-\tau_{i\to n})$, where
$\tau_{j\to n}$ is the mean first-passage time from $j$ to $n$
and $\tau_{n\to n}=0$~\cite{Meyer1975}.

\section{Discussion and conclusions}
\label{sec:discussion}

The argument applies to any linear
system $A\bm x=\bm b$ with nonsingular $A$ and fixed $\bm b$.
If only $A_{ij}$ changes and $A'$ is nonsingular, subtracting the
equations and eliminating $\Delta A_{ij}x_j'$ between $n$ and $m$ gives
\begin{align}
    \Delta x_n&=-(A^{-1})_{ni}\Delta A_{ij}x_j'\,,\nonumber\\
    \Delta x_n&=\frac{(A^{-1})_{ni}}{(A^{-1})_{mi}}\Delta x_m\,,
    \qquad (A^{-1})_{mi}\ne0\,.
    \label{eq:linear-system}
\end{align}
The same proportionality holds for fixed linear observables.

The same argument applies whenever $\Delta A\bm x'=\bm u\alpha$,
with fixed $\bm u$: the perturbation enters the equations in fixed
proportions, with only its scalar amplitude $\alpha$ changing.
Then $\Delta\bm x=-A^{-1}\bm u\alpha$, so eliminating $\alpha$
again gives mutual linearity.

The frequency-domain results follow by setting $A=sI-W$.
For time-independent rates, the Laplace-transformed master equation is
$(sI-W)\widetilde{\bm p}(s)=\bm p(0)$.
Keeping the same initial distribution in both systems, define
$R(s)=(sI-W)^{-1}$ for $\operatorname{Re}s>0$ and
$h_n(s)=R_{ni}(s)-R_{nj}(s)$.
Changing the rates of edge $(i,j)$ and eliminating the common factor
between states $n$ and $m$ gives
\begin{align}
    \Delta\widetilde p_n
    &=h_n(s)(\Delta W_{ij}\widetilde p_j'
              -\Delta W_{ji}\widetilde p_i')\,,\nonumber\\
    \Delta\widetilde p_n
    &=\frac{h_n(s)}{h_m(s)}\Delta\widetilde p_m\,,
    \qquad h_m(s)\ne0\,.
    \label{eq:laplace}
\end{align}
All transformed probabilities are evaluated at the same $s$.
This recovers the Laplace-domain probability
relations~\cite{ZhengLu2026,VoitsSchwarz2026}.
Repeating \cref{eq:output-current,eq:input-current,eq:currents} with
$\widetilde{\bm p}(s)$ and $h_n(s)$ gives the corresponding current
relations~\cite{Harunari2024}.
The frequency-dependent slope corresponds to a convolution in time,
rather than an affine relation at each instant.

\bibliography{references}
\end{document}
